\documentclass[11pt]{article}
\usepackage{amsmath,amssymb}

\title{A Small Polynomial}
\author{Feather demo}
\date{}

\begin{document}
\maketitle

\section{Evaluation}
Consider the polynomial $p(x) = 1 + 2x + 3x^2$ over the real numbers.
At $x = 2$, we have
\[
p(2) = 1 + 2 \cdot 2 + 3 \cdot 2^2 = 17.
\]

\section{Horner's method}
The same polynomial can be written as
\begin{align*}
p(x) &= 1 + x(2 + 3x), \\
p(2) &= 1 + 2(2 + 6) = 17.
\end{align*}
This form evaluates the polynomial without computing powers separately.

\section{A simple identity}
The sum of the first $n$ positive integers is
\[
\sum_{k=1}^{n} k = \frac{n(n+1)}{2}.
\]

\section{A derivative}
The derivative measures the slope of the graph:
\[
p'(x) = 2 + 6x, \qquad p'(2) = 14.
\]
The second derivative is $p''(x) = 6$, so the graph curves upward everywhere.

\section{An integral}
An antiderivative is
\[
\int p(x)\,dx = x + x^2 + x^3 + C.
\]
Since $p(x)$ is positive on $[0,1]$, the area under its graph on this interval is
\begin{align*}
\int_0^1 p(x)\,dx
  &= \left[x + x^2 + x^3\right]_0^1 \\
  &= (1 + 1 + 1) - 0 = 3.
\end{align*}

\section{Evaluation as a matrix}
Evaluating at three points can be written as one matrix multiplication:
\[
\begin{pmatrix}
p(0) \\ p(1) \\ p(2)
\end{pmatrix}
=
\begin{pmatrix}
1 & 0 & 0 \\
1 & 1 & 1 \\
1 & 2 & 4
\end{pmatrix}
\begin{pmatrix} 1 \\ 2 \\ 3 \end{pmatrix}
=
\begin{pmatrix} 1 \\ 6 \\ 17 \end{pmatrix}.
\]
Each row evaluates the same coefficient vector at a different input.

\end{document}
